\chapter{Derivación de las ecuaciones de Newton--Euler}\label{ch:eom}
\objetivos{Construir el modelo rígido desde los balances; conservar los signos de los productos de inercia; reducirlo a vuelo simétrico e identificar las hipótesis terrestres.}

\section{Cantidad de movimiento del centro de gravedad}
Se considera un avión rígido de masa constante durante el intervalo dinámico y un campo gravitatorio uniforme. En una referencia inercial,
\begin{equation}
\frac{\dd}{\dd t}(m\vect v_G)=\sum\vect F=m\vect a_G.
\end{equation}
La fuerza resultante contiene aerodinámica, propulsión y peso. Se define $(X,Y,Z)$ como fuerza no gravitatoria total en cuerpo. Al aplicar transporte a la velocidad y añadir la gravedad expresada en cuerpo,
\begin{align}
\dot u&=rv-qw+X/m-g\sin\theta,\\
\dot v&=pw-ru+Y/m+g\sin\phi\cos\theta,\\
\dot w&=qu-pv+Z/m+g\cos\phi\cos\theta.
\label{eq:traslacionales}
\end{align}
Estas velocidades son respecto al suelo inercializado. Si se quieren usar velocidades relativas al aire en un campo de viento variable, se resta la aceleración del viento descrita en el capítulo~\ref{ch:cinematica}. Las fuerzas aerodinámicas se calculan en cualquier caso con la velocidad aérea.

\subsection{Por qué no aparece un término arbitrario de masa variable}
En misiones se actualiza la masa por consumo, pero el modelo instantáneo de fuerzas del motor ya representa el intercambio de cantidad de movimiento con el flujo propulsivo. Añadir sin más $-\dot m\vect v$ a las ecuaciones del avión introduce dependencia del origen de velocidades y puede duplicar el empuje. Una derivación completa de sistema abierto requiere flujos de cantidad de movimiento y su velocidad relativa. La aproximación cuasiestacionaria utilizada aquí aplica las ecuaciones rígidas con la masa instantánea y contabiliza el consumo por separado. No pretende representar expulsión impulsiva de carga ni cambios rápidos de distribución de masa.

\section{Balance de momento angular}
Respecto al centro de gravedad, el momento angular es
\begin{equation}
\vect H_G^b=I\vect\omega^b.
\end{equation}
La ecuación vectorial inercial $\vect M_G=(\dd\vect H_G/\dd t)_i$ se expresa en cuerpo mediante transporte. Con tensor constante en cuerpo,
\begin{equation}
I\dot{\vect\omega}+\vect\omega\times(I\vect\omega)=\vect M_G.
\label{eq:euler-vector}
\end{equation}
Si el tensor varía, aparece además $\dot I\vect\omega$ y, según el mecanismo de variación, los flujos de momento angular correspondientes. Esa situación exige un modelo adicional y no debe ocultarse dentro de una inercia media.

\section{Expansión para un avión simétrico}
Se usa el convenio del capítulo~\ref{ch:fuerzas}: el producto integral $I_{xz}$ aparece con signo negativo en el tensor. El momento angular es
\begin{equation}
\vect H^b=(I_{xx}p-I_{xz}r,\ I_{yy}q,\ I_{zz}r-I_{xz}p)^T.
\end{equation}
Multiplicar por la matriz del producto vectorial, y no memorizar el resultado, proporciona
\begin{align}
\mathcal L&=I_{xx}\dot p-I_{xz}\dot r+(I_{zz}-I_{yy})qr-I_{xz}pq,\\
\mathcal M&=I_{yy}\dot q+(I_{xx}-I_{zz})pr+I_{xz}(p^2-r^2),\\
\mathcal N&=I_{zz}\dot r-I_{xz}\dot p+(I_{yy}-I_{xx})pq+I_{xz}qr.
\label{eq:euler-componentes}
\end{align}
Los términos cuadráticos son los acoplamientos giroscópicos del sólido. Aunque la aeronave sea geométricamente simétrica, $I_{xz}$ puede ser no nulo si los ejes elegidos no son principales.

Para integrar, se puede resolver directamente el sistema lineal con $I$ en cada paso. En el plano $p,r$ el determinante es $\Delta=I_{xx}I_{zz}-I_{xz}^2$. Definiendo
\begin{align}
A_p&=\mathcal L-(I_{zz}-I_{yy})qr+I_{xz}pq,\\
A_r&=\mathcal N-(I_{yy}-I_{xx})pq-I_{xz}qr,
\end{align}
se obtiene
\begin{equation}
\dot p=\frac{I_{zz}A_p+I_{xz}A_r}{\Delta},\qquad
\dot r=\frac{I_{xz}A_p+I_{xx}A_r}{\Delta}.
\end{equation}
Si $\Delta$ es muy pequeño, el problema está mal condicionado o la distribución de masa es casi degenerada. No procede dividir por él sin examinar el tensor.

\section{Comprobaciones internas de la dinámica}
\subsection{Límite de ejes principales}
Al imponer $I_{xz}=0$, desaparece el acoplamiento entre aceleraciones $\dot p$ y $\dot r$. Si además $I_{xx}=I_{yy}=I_{zz}=I_s$, se anula $\vect\omega\times I_s\vect\omega$ y resulta $I_s\dot{\vect\omega}=\vect M$. El caso de un sólido con inercia isotrópica es una prueba útil del álgebra.

\subsection{Energía de rotación}
La energía cinética angular es $E_r=\tfrac12\vect\omega^TI\vect\omega$. Para tensor constante,
\begin{equation}
\dot E_r=\vect\omega^TI\dot{\vect\omega}
=\vect\omega^T\vect M,
\end{equation}
porque $\vect\omega\cdot(\vect\omega\times I\vect\omega)=0$. Sin momento externo, la energía se conserva aunque las componentes de velocidad angular cambien. Esta prueba detecta varios errores de signo, pero no todos; se acompaña de la comprobación componente a componente.

\caso{Ejemplo: producto de inercia no nulo}{Considérese
\[
I=\begin{bmatrix}8500&0&-620\\0&13000&0\\-620&0&17000\end{bmatrix}\SIu{kg\,m^2},\quad
\vect\omega=(0.12,-0.08,0.15)\SIu{rad\,s^{-1}},
\]
y $\vect M=(950,-500,340)\SIu{N\,m}$. Se calcula primero $\vect H=I\vect\omega$ y después $\dot{\vect\omega}=I^{-1}[\vect M-\vect\omega\times\vect H]$. Sustituir el resultado en la ecuación vectorial recupera los tres momentos con residuo menor que $10^{-10}\SIu{N\,m}$ en doble precisión. El residuo comprueba esta operación, no la correspondencia de esos parámetros con un avión real. El archivo reproducible conserva el vector numérico sin redondeo.
\textbf{Momento angular, producto vectorial y sistema de ecuaciones.} Primero se calculan las tres filas de $\boldsymbol H=I\boldsymbol\omega$. Las entradas $xz$ y $zx$ son negativas con el convenio del enunciado.\[\fitmath{\boldsymbol H=\left[\begin{matrix}927.0\\-1040.0\\2475.6\end{matrix}\right]\,\mathrm{kg\,m^2/s}}\]\[\fitmath{H_{1}=(8500)(0.12)+(-620)(0.15)\approx927\,\mathrm{kg\,m^2/s}}\]\[\fitmath{H_{2}=(13000)(-0.08)\approx-1040\,\mathrm{kg\,m^2/s}}\]\[\fitmath{H_{3}=(-620)(0.12)+(17000)(0.15)\approx2475.6\,\mathrm{kg\,m^2/s}}\]\[\fitmath{c_{1}=(-0.08)(2475.6)-(0.15)(-1040)\approx-42.048\,\mathrm{N\,m}}\]\[\fitmath{c_{2}=(0.15)(927)-(0.12)(2475.6)\approx-158.022\,\mathrm{N\,m}}\]\[\fitmath{c_{3}=(0.12)(-1040)-(-0.08)(927)\approx-50.64\,\mathrm{N\,m}}\]El producto giroscópico se desarrolla como $(qH_z-rH_y,\ rH_x-pH_z,\ pH_y-qH_x)$; su valor es:\[\fitmath{\boldsymbol\omega\times\boldsymbol H=\left[\begin{matrix}-42.048\\-158.022\\-50.64\end{matrix}\right]\,\mathrm{N\,m}}\]Se resta de los momentos aplicados, obteniendo el segundo miembro $\boldsymbol b=\boldsymbol M-\boldsymbol\omega\times\boldsymbol H$.\[\fitmath{\boldsymbol b=\left[\begin{matrix}992.048\\-341.978\\390.64\end{matrix}\right]\,\mathrm{N\,m}}\]La ecuación de $y$ queda desacoplada: $\dot q=b_y/I_{yy}$. Para $x,z$, eliminar la otra aceleración en el sistema $I_{xx}\dot p-I_{xz}\dot r=b_x$, $-I_{xz}\dot p+I_{zz}\dot r=b_z$ da:\[\fitmath{\dot p=\frac{I_{zz}b_x+I_{xz}b_z}{I_{xx}I_{zz}-I_{xz}^2},\qquad\dot r=\frac{I_{xz}b_x+I_{xx}b_z}{I_{xx}I_{zz}-I_{xz}^2}}\]\[\fitmath{(\dot p,\dot q,\dot r)=\left[\begin{matrix}0.118703407542279 & -0.026306 & 0.0273080066280125\end{matrix}\right]\,\mathrm{s^{-2}}}\]Para comprobar, multiplicar el vector de aceleraciones por $I$ y sumar el término giroscópico recupera $\boldsymbol M$. El trabajo giroscópico se anula por perpendicularidad; la tasa de energía es $\boldsymbol\omega\cdot\boldsymbol M=205\,\mathrm W$.

}

\section{Cierre del modelo de seis grados de libertad}
Las tres ecuaciones traslacionales, tres rotacionales, tres de actitud y tres de posición forman un sistema de doce ecuaciones de primer orden. En forma compacta,
\begin{equation}
\dot{\vect x}=\vect f(\vect x,\vect\delta,\vect a,t),
\end{equation}
donde $\vect\delta$ contiene mandos y $\vect a$ el ambiente. Para cerrar el sistema hacen falta modelos de aerodinámica y propulsión, masa e inercia, atmósfera y gravedad. Cada estado inicial debe ser compatible con esos modelos. Un XML válido no garantiza esa compatibilidad.

Las ecuaciones descritas usan Tierra plana y NED aproximadamente inercial. JSBSim propaga sobre una Tierra rotatoria y contiene más términos. La comparación correcta consiste en identificar qué efectos se han despreciado y su escala; no en exigir igualdad exacta de cada salida con el modelo plano.

\section{Reducción a movimiento simétrico}
En vuelo simétrico con $v=p=r=\phi=0$, sin fuerza lateral ni momentos de alabeo y guiñada, se obtiene
\begin{align}
\dot u&=-qw+X/m-g\sin\theta,\\
\dot w&=qu+Z/m+g\cos\theta,\\
I_{yy}\dot q&=\mathcal M,\qquad \dot\theta=q.
\end{align}
No se ha impuesto todavía $q=0$ ni $\dot u=\dot w=0$. La reducción simétrica permite cabeceo y cambios de trayectoria. Con $u=V\cos\alpha$, $w=V\sin\alpha$, se rota la dinámica al eje de velocidad. Si el empuje forma $\varepsilon$ con la trayectoria y $\gamma=\theta-\alpha$,
\begin{align}
m\dot V&=T\cos\varepsilon-D-W\sin\gamma,\\
mV\dot\gamma&=L+T\sin\varepsilon-W\cos\gamma.
\label{eq:simetrico-trayectoria}
\end{align}
Al prescindir de la ecuación de momento se supone que el mando puede mantener la incidencia exigida. Para un problema de trim o de estabilidad longitudinal hay que recuperarla.

\section{Equilibrio rectilíneo y viraje estacionario}
En vuelo rectilíneo horizontal simétrico con empuje alineado con trayectoria, $T=D$, $L=W$ y $\mathcal M=0$. Las tres condiciones son necesarias. En un viraje, $\dot\psi\neq0$ y las componentes $(p,q,r)$ se calculan con la cinemática de Euler. No se anulan todas las velocidades angulares por conservar alabeo y cabeceo.

Para $\dot\phi=\dot\theta=0$,
\begin{equation}
p=-\dot\psi\sin\theta,\quad
q=\dot\psi\sin\phi\cos\theta,\quad
r=\dot\psi\cos\phi\cos\theta.
\end{equation}
Esta relación explica por qué un trim completo de viraje es distinto de equilibrar sólo sustentación y peso.

\verifica{Se comprobó simbólicamente la expansión del tensor con $-I_{xz}$ y la anulación del trabajo giroscópico. El residuo vectorial se evalúa además en un caso con producto cruzado no nulo. Las hipótesis de rigidez, distribución de masa y fuerzas permanecen sujetas a revisión física.}
\section{Demostración desde un continuo de partículas}
\subsection{Separar traslación y rotación}
Considérese un conjunto material de partículas y un origen G situado en su centro de masa. Para cada partícula, $\vect r_i=\vect r_G+\vect\rho_i$ y $\sum m_i\vect\rho_i=0$. La cantidad de movimiento total es
\begin{equation}
\vect P=\sum_i m_i\dot{\vect r}_i
=m\dot{\vect r}_G+\frac{\dd}{\dd t}\sum_i m_i\vect\rho_i
=m\vect v_G,
\end{equation}
si las masas de las partículas permanecen constantes. Al sumar sus ecuaciones de Newton, las fuerzas internas se cancelan por pares bajo las hipótesis de interacción que permiten modelar el sólido clásico. Así se obtiene $m\vect a_G=\sum\vect F_{ext}$ sin atribuir al centro de masa una estructura interna ficticia.

El momento angular respecto a G se define con velocidades relativas a G:
\begin{equation}
\vect H_G=\sum_i\vect\rho_i\times m_i\dot{\vect\rho}_i.
\end{equation}
Para un sólido rígido, $\dot{\vect\rho}_i=\vect\omega\times\vect\rho_i$. La identidad de triple producto da
\begin{equation}
\vect\rho_i\times(\vect\omega\times\vect\rho_i)
=[(\vect\rho_i^T\vect\rho_i)I_3-\vect\rho_i\vect\rho_i^T]\vect\omega.
\end{equation}
Sumando se obtiene el tensor del capítulo anterior y $\vect H_G=I\vect\omega$. Esta demostración explica el signo negativo de los elementos no diagonales: procede del término $-\vect\rho\vect\rho^T$, no de una elección posterior para ajustar ecuaciones.

\subsection{Por qué G simplifica el balance de momentos}
Al derivar el momento angular alrededor de un punto móvil aparecen términos asociados al movimiento del origen. En G, las contribuciones que contienen $\sum m_i\vect\rho_i$ se anulan. Por ello el momento externo respecto a G se iguala a la derivada inercial de $\vect H_G$. Elegir otro punto fijo al avión exige conservar los términos de traslación del origen; no basta con reutilizar la misma ecuación con un tensor trasladado.

La energía cinética total se separa de forma análoga:
\begin{equation}
E_k=\frac12m\vect v_G^T\vect v_G+\frac12\vect\omega^TI\vect\omega.
\end{equation}
El término cruzado desaparece porque G es centro de masa. En un campo uniforme se puede añadir $mgh_G$. La tasa de energía se compara con trabajo de fuerzas resultantes y momentos en un sistema cerrado. Para un avión con flujo propulsivo y consumo, esa identidad exige el tratamiento de sistema abierto si se interpreta como balance de toda la energía del combustible y la estela.

\section{Tierra rotatoria: derivación de los términos adicionales}
Sea una referencia e que gira uniformemente respecto a una inercial i con velocidad angular $\vect\Omega$. Para una posición $\vect r$ respecto al centro terrestre,
\begin{equation}
\vect v_i=\vect v_e+\vect\Omega\times\vect r.
\end{equation}
Derivando con transporte y $\dot{\vect\Omega}=0$,
\begin{equation}
\vect a_i=\vect a_e+2\vect\Omega\times\vect v_e
+\vect\Omega\times(\vect\Omega\times\vect r).
\end{equation}
La ecuación en e resulta
\begin{equation}
\vect a_e=\frac{\vect F_{no\,g}}m+\vect g_{grav}
-2\vect\Omega\times\vect v_e
-\vect\Omega\times(\vect\Omega\times\vect r).
\end{equation}
Si se define gravedad efectiva como $\vect g_{ef}=\vect g_{grav}-\vect\Omega\times(\vect\Omega\times\vect r)$, la centrífuga ya está incluida en ella. Si la rotación no fuese uniforme, se añadiría el término de Euler $-\dot{\vect\Omega}\times\vect r$.

Una base NED que se desplaza introduce además velocidad de transporte del horizonte. No es la misma velocidad angular que el giro terrestre. Para un modelo de navegación global se distinguen tasa cuerpo respecto a inercial, tasa cuerpo respecto a Tierra y tasa horizonte respecto a Tierra. Las ecuaciones simplificadas del manual se recuperan al despreciar esas rotaciones y curvatura en una región local.

\verifica{El orden de magnitud de Coriolis es $2\Omega V$. Para una velocidad de unos $100\SIu{m/s}$ y una rotación terrestre de orden $7.3\times10^{-5}\SIu{s^{-1}}$, la escala es aproximadamente $0.015\SIu{m/s^2}$. Puede ser pequeña frente a g y a la vez importante para comparar residuos de trim de $10^{-4}\SIu{m/s^2}$. ``Pequeño para actuaciones'' no equivale a ``nulo en una auditoría de simulación''.}

\section{Pasar del modelo completo al balance de trayectoria}
Para $\beta=0$, $u=V\cos\alpha$ y $w=V\sin\alpha$. Derivando,
\begin{align}
\dot u&=\dot V\cos\alpha-V\dot\alpha\sin\alpha,\\
\dot w&=\dot V\sin\alpha+V\dot\alpha\cos\alpha.
\end{align}
Multiplicar la ecuación de u por $\cos\alpha$, la de w por $\sin\alpha$ y sumar cancela los términos de $\dot\alpha$ y los productos de q, obteniendo la ecuación tangencial. Para la ecuación normal se toma $-\sin\alpha$ por la primera y $\cos\alpha$ por la segunda. Aparece $V(\dot\alpha-q)$; con $\gamma=\theta-\alpha$ y $q=\dot\theta$, se convierte en $-V\dot\gamma$. Los signos de peso y sustentación completan la ecuación~\eqref{eq:simetrico-trayectoria}.

Esta reducción muestra de dónde procede el término de curvatura $mV\dot\gamma$. No es una corrección opcional para maniobras: es la aceleración normal del vector velocidad. Al imponer trayectoria rectilínea se anula porque $\dot\gamma=0$, no porque se haya supuesto velocidad de módulo constante.


\section{Taller resuelto: Dinámica y trabajo giroscópico}

\noindent\textbf{Nivel:} Intermedio.

\apartado{1. Enunciado}{Para el tensor Ixx=9800, Iyy=14800, Izz=19200 e Ixz=760 kg m², omega=(0.13,0.05,-0.09) rad/s y momentos (600,-380,420) N m, calcule aceleraciones y tasa de energía de rotación.}

\apartado{2. Datos}{El tensor usa -Ixz en entradas xz y zx y es constante en cuerpo.}

\apartado{3. Incógnitas}{Las magnitudes y comprobaciones solicitadas en el caso.}

\apartado{4. Sistema de referencia}{Cuerpo y momentos respecto al CG; producto integral Ixz positivo según dato.}

\apartado{5. Hipótesis físicas}{Sólido rígido y masa constante en el intervalo dinámico.}

\apartado{6. Modelo matemático}{Dinámica y trabajo giroscópico}

\apartado{7. Ecuaciones simbólicas}{Se parte de los balances y definiciones siguientes.}

\begin{equation}I\dot{\vect\omega}=\vect M-\vect\omega\times I\vect\omega\end{equation}

\begin{equation}\dot E_r=\vect\omega^TI\dot{\vect\omega}=\vect\omega^T\vect M\end{equation}

\apartado{8. Desarrollo algebraico}{Las operaciones se realizan en este orden, conservando las condiciones del modelo.}

\begin{enumerate}

\item Se calcula el momento angular antes del término giroscópico.

\item Se resuelve el sistema con el tensor completo.

\item Se evalúa la tasa de energía por dos vías: omega por M y omega por I odot.

\item El producto giroscópico no aporta trabajo porque es perpendicular a omega.

\end{enumerate}

\apartado{9. Sustitución numérica}{Se conservan los resultados intermedios sin redondear y se muestran todas las operaciones escalares disponibles.}

\textbf{Momento angular, producto vectorial y sistema de ecuaciones.} Primero se calculan las tres filas de $\boldsymbol H=I\boldsymbol\omega$. Las entradas $xz$ y $zx$ son negativas con el convenio del enunciado.\[\fitmath{\boldsymbol H=\left[\begin{matrix}1342.4\\740.0\\-1826.8\end{matrix}\right]\,\mathrm{kg\,m^2/s}}\]\[\fitmath{H_{1}=(9800)(0.13)+(-760)(-0.09)\approx1342.4\,\mathrm{kg\,m^2/s}}\]\[\fitmath{H_{2}=(14800)(0.05)\approx740\,\mathrm{kg\,m^2/s}}\]\[\fitmath{H_{3}=(-760)(0.13)+(19200)(-0.09)\approx-1826.8\,\mathrm{kg\,m^2/s}}\]\[\fitmath{c_{1}=(0.05)(-1826.8)-(-0.09)(740)\approx-24.74\,\mathrm{N\,m}}\]\[\fitmath{c_{2}=(-0.09)(1342.4)-(0.13)(-1826.8)\approx116.668\,\mathrm{N\,m}}\]\[\fitmath{c_{3}=(0.13)(740)-(0.05)(1342.4)\approx29.08\,\mathrm{N\,m}}\]El producto giroscópico se desarrolla como $(qH_z-rH_y,\ rH_x-pH_z,\ pH_y-qH_x)$; su valor es:\[\fitmath{\boldsymbol\omega\times\boldsymbol H=\left[\begin{matrix}-24.74\\116.668\\29.08\end{matrix}\right]\,\mathrm{N\,m}}\]Se resta de los momentos aplicados, obteniendo el segundo miembro $\boldsymbol b=\boldsymbol M-\boldsymbol\omega\times\boldsymbol H$.\[\fitmath{\boldsymbol b=\left[\begin{matrix}624.74\\-496.668\\390.92\end{matrix}\right]\,\mathrm{N\,m}}\]La ecuación de $y$ queda desacoplada: $\dot q=b_y/I_{yy}$. Para $x,z$, eliminar la otra aceleración en el sistema $I_{xx}\dot p-I_{xz}\dot r=b_x$, $-I_{xz}\dot p+I_{zz}\dot r=b_z$ da:\[\fitmath{\dot p=\frac{I_{zz}b_x+I_{xz}b_z}{I_{xx}I_{zz}-I_{xz}^2},\qquad\dot r=\frac{I_{xz}b_x+I_{xx}b_z}{I_{xx}I_{zz}-I_{xz}^2}}\]\[\fitmath{(\dot p,\dot q,\dot r)=\left[\begin{matrix}0.0655291072083522 & -0.0335586486486487 & 0.0229542771603306\end{matrix}\right]\,\mathrm{s^{-2}}}\]Para comprobar, multiplicar el vector de aceleraciones por $I$ y sumar el término giroscópico recupera $\boldsymbol M$. El trabajo giroscópico se anula por perpendicularidad; la tasa de energía es $\boldsymbol\omega\cdot\boldsymbol M=21.2\,\mathrm W$.





\apartado{10. Resultados con unidades}{Los resultados pertenecen al modelo y al intervalo de datos declarados.}

\begin{center}\begin{tabular}{lr}\toprule Magnitud&Resultado\\\midrule

pdot & $0.06552911\,\mathrm{s^{-2}}$\\

qdot & $-0.03355865\,\mathrm{s^{-2}}$\\

rdot & $0.02295428\,\mathrm{s^{-2}}$\\

Tasa de energía & $21.2\,\mathrm W$\\

\bottomrule\end{tabular}\end{center}

\apartado{11. Interpretación}{La energía puede aumentar o disminuir según omega por M, aun cuando algunas tasas angulares tengan signo distinto.}

\apartado{12. Comprobación independiente y alcance}{Residuo vectorial y igualdad de las dos tasas energéticas, ejecutados con álgebra matricial.}

\apartado{13. Limitaciones}{No se interpreta un literal de aceleración como prueba del balance; las operaciones y residuos están conservados en el programa.}

\par\smallskip\noindent{\footnotesize\textit{Datos y solución originales. Las sustituciones registradas se reproducen con SymPy y Studium; eso no establece validación experimental ni revisión experta.}}


\repaso{¿Permite la simetría geométrica imponer siempre $I_{xz}=0$?}{No. La simetría en $xz$ anula $I_{xy}$ e $I_{yz}$. Para anular $I_{xz}$ hay que elegir ejes principales en ese plano o disponer de una distribución de masa que lo haga cero.}
\resumen{Los balances se derivan en un sistema inercial y se expresan en cuerpo mediante transporte. Los signos de los productos de inercia dependen de su definición integral. La reducción de punto material necesita hipótesis de trim y no sustituye el equilibrio de momentos.}
\fuentes{How, MIT 16.333, lección 4, pp. 4--2 a 4--3. Esa fuente usa un elemento cruzado con otra notación; aquí se define explícitamente la integral y se rehace la expansión. JSBSim 1.3.3: FGAccelerations, FGPropagate y FGMassBalance como referencia de implementación, no como identidad del modelo terrestre simplificado.}
